% =========================================================
% Definition: Ring
% Section: Rings of Sets — Volume I
% Source: Halmos & Givant, Introduction to Boolean Algebras, §1
% =========================================================

\begin{tcolorbox}[colback=gray!6, colframe=gray!40, arc=2pt,
  left=6pt, right=6pt, top=4pt, bottom=4pt,
  title={\small\textbf{Toolkit --- Rings of Sets}},
  fonttitle=\small\bfseries]
\textbf{Concept family:} The ring axioms (Halmos--Givant §1) as the algebraic
foundation for Boolean algebras of sets.

\medskip
\textbf{Atomic items in this section:}
\begin{itemize}
  \item Definition: Ring
\end{itemize}

\medskip
\textbf{Key predicate:}
\[
  \operatorname{Ring}(R,+,\cdot,-,0,1)
  \;\colonequiv\;
  \operatorname{AbelianGroup}(R,+,-,0)
  \;\land\;
  \operatorname{Monoid}(R,\cdot,1)
  \;\land\;
  \operatorname{Distributes}(\cdot,+,R).
\]
\textbf{Axiom numbering} follows Halmos--Givant; axiom (4) (commutativity of
multiplication) is not a ring axiom but a theorem in the Boolean setting.
\end{tcolorbox}

\begin{tcolorbox}[colback=defbox, colframe=defborder, arc=2pt,
  left=6pt, right=6pt, top=4pt, bottom=4pt,
  title={\small\textbf{Definition --- Ring}},
  fonttitle=\small\bfseries]
\begin{definition}[Ring]
\label{def:ring}
A \textbf{ring} is a non-empty set $R$ together with two binary
operations $+$ and $\cdot$ on $R$, a unary operation $-$ on $R$,
and two distinguished elements $0, 1 \in R$, satisfying the
following axioms for all $p, q, r \in R$:
\begin{enumerate}
  \item[(1)] $p + (q + r) = (p + q) + r$
             \hfill\textit{(additive associativity)}
  \item[(2)] $p \cdot (q \cdot r) = (p \cdot q) \cdot r$
             \hfill\textit{(multiplicative associativity)}
  \item[(3)] $p + q = q + p$
             \hfill\textit{(additive commutativity)}
  \item[(5)] $p + 0 = p$
             \hfill\textit{(additive identity)}
  \item[(6)] $p \cdot 1 = p$
             \hfill\textit{(multiplicative identity)}
  \item[(7)] $p + (-p) = 0$
             \hfill\textit{(additive inverse)}
  \item[(8)] $p \cdot (q + r) = p \cdot q + p \cdot r$
             \hfill\textit{(left distributivity)}
  \item[(9)] $(q + r) \cdot p = q \cdot p + r \cdot p$
             \hfill\textit{(right distributivity)}
\end{enumerate}
Axiom numbering follows Halmos--Givant~\cite{HalmosGivant2009};
axiom~(4) is commutativity of multiplication, which is not a ring
axiom but a theorem in the Boolean setting.
A ring possessing a distinguished element $1$ satisfying
axiom~(6) is called a \textbf{ring with unit}.
A ring satisfying $p \cdot q = q \cdot p$ for all $p, q \in R$
is called a \textbf{commutative ring}.
\end{definition}
\end{tcolorbox}

\begin{remark*}[Standard quantified statement]
\[
\begin{aligned}
(R,+,\cdot,-,0,1)\ \text{is a ring}
\;\Longleftrightarrow\;
  &\forall p,q,r \in R{:} \\
  &\quad p + (q + r) = (p + q) + r \\
  &\;\land\quad p \cdot (q \cdot r) = (p \cdot q) \cdot r \\
  &\;\land\quad p + q = q + p \\
  &\;\land\quad p + 0 = p \\
  &\;\land\quad p \cdot 1 = p \\
  &\;\land\quad p + (-p) = 0 \\
  &\;\land\quad p \cdot (q + r) = p \cdot q + p \cdot r \\
  &\;\land\quad (q + r) \cdot p = q \cdot p + r \cdot p.
\end{aligned}
\]
\end{remark*}

\begin{remark*}[Definition predicate reading]
\[
\operatorname{Ring}(R,+,\cdot,-,0,1)
\;\Longleftrightarrow\;
\operatorname{AbelianGroup}(R,+,-,0)
\;\land\;
\operatorname{Monoid}(R,\cdot,1)
\;\land\;
\operatorname{Distributes}(\cdot,+,R)
\]
\end{remark*}

\begin{remark*}[Negated quantified statement]
\[
\begin{aligned}
(R,+,\cdot,-,0,1)\ \text{is \textbf{not} a ring}
\;\Longleftrightarrow\;
  &\exists p,q,r \in R{:} \\
  &\quad p + (q + r) \neq (p + q) + r \\
  &\;\lor\quad p \cdot (q \cdot r) \neq (p \cdot q) \cdot r \\
  &\;\lor\quad p + q \neq q + p \\
  &\;\lor\quad p + 0 \neq p \\
  &\;\lor\quad p \cdot 1 \neq p \\
  &\;\lor\quad p + (-p) \neq 0 \\
  &\;\lor\quad p \cdot (q + r) \neq p \cdot q + p \cdot r \\
  &\;\lor\quad (q + r) \cdot p \neq q \cdot p + r \cdot p.
\end{aligned}
\]
\end{remark*}

\begin{remark*}[Failure mode decomposition]
The eight axioms partition into three structural groups:
\[
\underbrace{(1),(3),(5),(7)}_{\displaystyle\operatorname{AbelianGroup}(R,+,-,0)}
\qquad
\underbrace{(2),(6)}_{\displaystyle\operatorname{Monoid}(R,\cdot,1)}
\qquad
\underbrace{(8),(9)}_{\displaystyle\operatorname{Distributes}(\cdot,+,R)}
\]
A structure violating any axiom in the first group lacks a coherent
additive group.
A structure violating any axiom in the second group lacks a
multiplicative identity or associativity.
A structure violating the third group possesses two individually
well-behaved operations that do not interact coherently;
distributivity is the sole structural bridge between $+$ and
$\cdot$, and its failure is sufficient to destroy the ring property.
\end{remark*}

\begin{remark*}[Interpretation]
A ring abstracts the arithmetic of the integers.
The additive structure is always a commutative group.
The multiplicative structure is a monoid --- associative with
unit --- but need not be commutative; axiom~(4) of
Halmos--Givant, commutativity of multiplication, is absent from
the ring axioms.
In the Boolean setting, multiplicative commutativity is not
assumed but derived as a theorem from idempotence alone.
Distributivity encodes the coherent interaction between the
two operations: it is the axiom that makes $\cdot$ and $+$
form a single algebraic system rather than two independent ones.
\end{remark*}

\NoLocalDependencies
